\section{Notation and exact layer specifications}
\label{app:v4-layers}
\begin{longtable}{>{\raggedright\arraybackslash}p{0.21\textwidth}>{\raggedright\arraybackslash}p{0.72\textwidth}}
\toprule
Symbol & Meaning\\
\midrule
\endhead
$\mathbf{Y}^{\mathrm{tr}}$, $\vD$ & Binary target-training association matrix and its unique positive edges.\\
$\mathbf{b}_e,\mathbf{b}_r$ & Native 512-dimensional phase-1 protein and reaction outputs.\\
$\mathbf{g}_e,\mathbf{f}_e$ & Separately normalized global and fused protein branches.\\
$K$, $V=K+2$ & Learned residue-view count (default 4) and total protein-view count.\\
$i\in\mathcal V_e$, $k$, $v$, $d$ & Residue index; learned query index $1\le k\le K$; view label $v\in\mathcal I=\{P,S,1,\ldots,K\}$; feature coordinate $1\le d\le256$.\\
$\mathbf h_i$, $\ell_i$, $\tilde\ell_i$ & Frozen 1024-dimensional residue feature, SLEEC logit, and centered/clipped logit; enzyme index suppressed.\\
$\mathbf q_k$, $\hat{\mathbf q}_k$ & Shared trainable query and its numerically normalized version, both 256-dimensional.\\
$\mathbf k_i$, $\mathbf v_i$ & Normalized residue key and unnormalized value, both 256-dimensional; enzyme index suppressed.\\
$\mathbf x_e^{v}$ & Raw views: $v=P$ global ProtT5 and $v=S$ SLEEC (1024 dimensions); $v=k$ learned query $k$ (256 dimensions).\\
$\mathbf y_e^{v}$, $\bar{\mathbf y}_e$ & Projected, normalized views and their coordinate-wise fused vector, all 256-dimensional.\\
$p_{ki},a_{ki}$ & Raw and uniform-smoothed residue attention for learned slot $k$.\\
$\omega_{e,vd}$ & Gate weight for protein view $v$ and hidden coordinate $d$.\\
$s$ & Learned bounded protein fusion scale; initial value 0.1.\\
$\nu$ & Protein fusion multiplier used after training; value 3.\\
$h_E,h_R$ & Independent residual heads fitted with frozen F3 inputs.\\
$\kappa,c$ & Residual coefficient in training (0.2) and inference multiplier (0.5).\\
$\mathbf{u}_e^{\mathrm{tr}},\mathbf{u}_r^{\mathrm{tr}}$ & Refined endpoints used in the phase-2 training objectives.\\
$\mathbf{u}_e,\mathbf{u}_r$ & Refined endpoints after the two inference adjustments.\\
$s_{ax}^{(1)},s_{ax}^{(2)}$ & Scalar association logits in phases 1 and 2; distinct from the fusion scale $s$ and final embeddings $\mathbf z_x$.\\
$\hat{\mathbf u}_i$ & Renormalized phase-2 vector in the biological-supervision ablation only.\\
$\xi_c,w_{ic}$ & Category weight and membership confidence in the biological ablation.\\
$\mathbf{a}_e,\mathbf{a}_r$ & Normalized training-reaction dictionary coordinates.\\
$\mathbf{z}_e,\mathbf{z}_r$ & Complete $512+N_R$ endpoint representations.\\
$\alpha$ & Dictionary contribution to the final dot product; value 0.4.\\
$\mathcal M_r$, $\mathcal P_r$ & Available reaction modalities and multiset of valid participant molecules, respectively.\\
\bottomrule
\end{longtable}

\begin{longtable}{>{\raggedright\arraybackslash}p{0.30\textwidth}>{\raggedright\arraybackslash}p{0.63\textwidth}}
\caption{Active neural components in the default $K=4$ configuration.
Dimensions are listed in execution order.}\label{tab:v4-layers}\\
\toprule
Component & Specification\\
\midrule
\endfirsthead
\toprule Component & Specification\\\midrule
\endhead
Frozen residue features & ProtT5; 1024 coordinates per retained residue.\\
Frozen SLEEC scorer & Linear $1024\to256$; ReLU; linear $256\to1$.\\
Residue adapter & LayerNorm(1024); linear $1024\to256$; GELU.\\
Keys and values & Keys: bias-free linear $256\to256$ followed by normalization. Values: affine linear $256\to256$.\\
Learned residue queries & Four trainable 256-dimensional vectors; orthogonal initialization; normalized before scoring.\\
Global/SLEEC view projection & Two independent networks: LayerNorm(1024); linear $1024\to512$; GELU; dropout 0.1; linear $512\to256$; normalization.\\
Learned-slot view projection & Four independent networks: LayerNorm(256); linear $256\to512$; GELU; dropout 0.1; linear $512\to256$; normalization.\\
Protein gate & Linear $1536\to256$; GELU; linear $256\to1536$; reshape to $6\times256$; softmax over views.\\
Global output branch & LayerNorm(1024); linear $1024\to512$; GELU; dropout 0.1; linear $512\to512$; normalization.\\
Fused output branch & LayerNorm(256); linear $256\to512$; GELU; dropout 0.1; linear $512\to512$; normalization.\\
Uni-Mol2 molecule pooling & Separate scalar affine attention for reactants and products, acting on 768-dimensional molecule vectors.\\
ChIRo molecule pooling & Separate scalar affine attention for reactants and products, acting on 256-dimensional molecule vectors.\\
Reaction modality MLPs & Input widths $768,3072,1024,617$ for T5, Uni-Mol2 composition, ChIRo composition, and chemistry. Each follows $d_m\to1024\to512\to512$; ReLU after its first two layers.\\
Reaction token normalization & One LayerNorm(512), shared over modality tokens. No per-token $\ell_2$ normalization in this configuration.\\
Reaction modality attention & Linear $512\to512$; ReLU; linear $512\to1$; masked softmax over available modalities.\\
Reaction output projection & Linear $512\to1024$; ReLU; linear $1024\to512$; normalization.\\
Phase-2 heads & Two independent copies of LayerNorm(512); linear $512\to1024$; GELU; linear $1024\to512$. Last affine layer initialized to zero.\\
\bottomrule
\end{longtable}

The two phase-2 heads contain 2,102,272 parameter values in total.
The four residue queries are shared across proteins; they are not four
trainable vectors created separately for each sequence.

\section{Attention implementation and sequence preprocessing}
\label{app:v4-attention}
Padding is removed from softmax denominators and pooling sums, and its
contribution is zero. Before the frozen scorer and residue adapter, padded
values are replaced by zero so that padded NaN sentinels do not enter learned
linear maps. SLEEC logits are detached. Centering its logits alone would not
alter softmax, but centering before clipping gives a protein-relative clipping
rule. The use of logits, rather than sigmoid scores, is material to the
resulting distribution.

The learned attention scale is initialized by
$\rho_\gamma=\operatorname{logit}(4/10)$, and the fusion scalar by
$\rho_s=\operatorname{logit}(0.1/0.5)$. The normalization floor in the
protein multiview module is $10^{-6}$. Attention entropy uses natural
logarithms and a probability floor of $10^{-12}$ inside the logarithm.
Regularization in Eqs.~\eqref{eq:v4-entropy}--\eqref{eq:v4-diversity} is
evaluated on raw learned attention, not on the SLEEC distribution and not
on the smoothed distributions used to pool values.

For the F3 protein tower, at most 1022 residue vectors are retained. For a
longer cached sequence, the deterministic \texttt{ends\_center} policy keeps
$\lfloor1022/4\rfloor$ vectors at each end and the remaining budget from
the center, preserving their original order. The dictionary's enzyme means
are exported separately by averaging the stored residue bank; they should
not be silently replaced with the neural tower's potentially truncated mean.
EnzymeMap's specified screening library already restricts sequence length
to at most 650 amino acids.

The pooling operations scale linearly in sequence length for a fixed number
of slots. In particular, learned residue attention forms an $L\times K$
matrix ($L\times4$ by default), not an $L\times L$ self-attention matrix. This complexity statement
concerns the trainable pooling module after frozen features are available;
it does not describe the cost of computing ProtT5 features.

\section{Fixed chemistry descriptor}
\label{app:v4-chemistry}
The fixed chemistry descriptor gives the reaction encoder an explicit summary
of participant substructures, physicochemical properties, and recognized
cofactors alongside its pretrained reaction and molecular embeddings. These
quantities are available directly from participant SMILES; their construction
does not require enzyme--reaction association labels. This supplies observable
chemical cues whose contribution to retrieval is learned through the modality
projection and fusion.

The 617-dimensional chemistry vector consists of four blocks:
\begin{equation}
  617=512+64+4+37.
\end{equation}
The first block averages 512-bit radius-2 Morgan fingerprints of valid
participant molecules. The second block applies sum, mean, maximum, and
standard-deviation aggregation to 16 molecule descriptors, yielding 64
coordinates. Table~\ref{tab:v4-molecular-properties} lists the properties in
their exact order within the per-molecule vector $\mathbf{d}(m)$. Each is computed
with RDKit from a valid participant molecule. The aggregate block concatenates
all 16 sums, then all 16 means, maxima, and population standard deviations;
these are summaries across participants, irrespective of reaction side.

\begin{longtable}{@{}r>{\raggedright\arraybackslash}p{0.26\textwidth}>{\raggedright\arraybackslash}p{0.62\textwidth}@{}}
\caption{The 16 molecular properties aggregated in the reaction chemistry
descriptor. Each property contributes four aggregate coordinates, giving
$16\times4=64$ in total. Function names identify the RDKit calculations.}
\label{tab:v4-molecular-properties}\\
\toprule
Index & Property & Definition or calculation\\
\midrule
\endfirsthead
\toprule
Index & Property & Definition or calculation (continued)\\
\midrule
\endhead
1 & Exact molecular weight & Exact molecular weight from \texttt{CalcExactMolWt}.\\
2 & Heavy-atom count & Number of atoms with atomic number $Z>1$.\\
3 & Heteroatom count & Number of atoms with $Z\notin\{0,1,6\}$; wildcard atoms, hydrogen, and carbon are excluded.\\
4 & Formal charge & Net molecular formal charge from \texttt{GetFormalCharge}.\\
5 & Hydrogen-bond donors & Donor count from \texttt{NumHDonors}.\\
6 & Hydrogen-bond acceptors & Acceptor count from \texttt{NumHAcceptors}.\\
7 & Topological polar surface area & Molecular TPSA from \texttt{CalcTPSA}.\\
8 & $\log P$ & RDKit Crippen estimate from \texttt{MolLogP}.\\
9 & Rotatable-bond count & Number of rotatable bonds from \texttt{NumRotatableBonds}.\\
10 & Ring count & Number of rings from \texttt{RingCount}.\\
11 & Aromatic-heavy-atom fraction & Number of aromatic heavy atoms divided by $\max(N_{\mathrm{heavy}},1)$, where $N_{\mathrm{heavy}}$ is the heavy-atom count.\\
12 & Atom stereocenter count & Number of centers returned by \texttt{FindMolChiralCenters}, including unassigned centers.\\
13 & Wildcard-atom count & Number of atoms with symbol \texttt{*}.\\
14 & Nitrogen count & Number of atoms with symbol \texttt{N}.\\
15 & Oxygen count & Number of atoms with symbol \texttt{O}.\\
16 & Phosphorus/sulfur count & Combined number of atoms with symbols \texttt{P} or \texttt{S}.\\
\bottomrule
\end{longtable}

The four set statistics are participant count, unique-participant fraction,
valid-participant fraction, and wildcard-containing fraction among valid
participants. The final block is a 37-position capacity for the recognized
core-cofactor vocabulary fitted on training reactions; unused positions are
zero. Only the 64 aggregated property coordinates are standardized. Their
means and standard deviations are fitted on valid training reactions, with a
standard deviation below $10^{-6}$ replaced by one. Fingerprints, set
statistics, and cofactor indicators are not standardized in this construction.

More precisely, let $\mathcal P_r$ be the multiset of valid participant
molecules, with $n_r=|\mathcal P_r|>0$. For a molecule $m$, let
$\boldsymbol\phi(m)\in\{0,1\}^{512}$ denote its radius-2 Morgan fingerprint and
$\mathbf{d}(m)\in\mathbb R^{16}$ its property vector. Define
\begin{align}
 \bar{\boldsymbol\phi}_r&=\frac{1}{n_r}\sum_{m\in\mathcal P_r}\boldsymbol\phi(m),\\
 \mathbf{d}_r&=\left[\sum_{m\in\mathcal P_r}\mathbf{d}(m);\;
       \operatorname{mean}_{m\in\mathcal P_r}\mathbf{d}(m);\;
       \operatorname{max}_{m\in\mathcal P_r}\mathbf{d}(m);\;
       \operatorname{std}_{m\in\mathcal P_r}\mathbf{d}(m)\right],\\
 \mathbf{c}_r^{\mathrm{chem}}&=
   [\bar{\boldsymbol\phi}_r;\;(\mathbf{d}_r-\boldsymbol\mu_{\mathrm{tr}})\oslash\boldsymbol\sigma_{\mathrm{tr}};
     \;\mathbf{s}_r^{\mathrm{set}};\;\mathbf{q}_r^{\mathrm{cof}}]\in\mathbb R^{617}.
 \label{eq:v4-chemistry-vector}
\end{align}
Here $\boldsymbol\mu_{\mathrm{tr}},\boldsymbol\sigma_{\mathrm{tr}}\in\mathbb R^{64}$ are the
training means and stabilized standard deviations defined above, and
$\oslash$ denotes elementwise division. All property aggregations and standardization are coordinate-wise, and
$\operatorname{std}$ uses the population convention. The vector
$\mathbf{s}_r^{\mathrm{set}}\in\mathbb R^4$ contains the four set statistics above;
$\mathbf{q}_r^{\mathrm{cof}}\in\{0,1\}^{37}$ records presence of recognized core
cofactors in the training vocabulary. An unseen cofactor does not create a
new coordinate. Reactions with no valid participant have zero fingerprint
and property blocks and an unavailable-chemistry mask; modality attention
excludes this input.

``Fixed'' therefore distinguishes feature construction from neural learning.
The fingerprint and property definitions, training-derived normalization,
and cofactor vocabulary remain unchanged while the reaction encoder is
trained. A learned chemistry-specific MLP maps $\mathbf{c}_r^{\mathrm{chem}}$ to a
512-dimensional token, which participates in the same modality attention as
the other three sources (Section~\ref{sec:v4-reaction}).

Participant canonicalization removes atom-map indices while retaining
isomeric SMILES information. When both sides of a stored self-reaction have
the same canonical participant multiset, the participant list is used once.
Otherwise, the two sides are combined for this descriptor. Thus this chemistry
block is a participant-set feature, even when the neural molecular branches
retain physical reaction sides. Cofactor-related participant features remain
part of the reaction representation; the separate annotation loss in
Section~\ref{sec:v4-biology} is used only in the biological ablation.

\section{Efficient evaluation of the biological ablation}
\label{app:v4-efficient-bio}
This section concerns only the auxiliary-loss ablation in
Section~\ref{sec:v4-biology}, not the proposed training objective.
For one category, define the weighted sums
\begin{equation}
  A_c=\sum_{i\in I_c}w_{ic},\qquad
  Q_c=\sum_{i\in I_c}w_{ic}^{2},\qquad
  \mathbf{s}_c=\sum_{i\in I_c}w_{ic}\hat{\mathbf{u}}_i,
\end{equation}
and the diagonal correction
$V_c=\sum_{i\in I_c}w_{ic}^{2}\|\hat{\mathbf{u}}_i\|_2^2$.
The ordered off-diagonal confidence mass is $T_c=A_c^2-Q_c$. Expanding the
dot product of the weighted sum gives
\begin{equation}
  \sum_{\substack{i,j\in I_c\\i\ne j}}
      w_{ic}w_{jc}(1-\hat{\mathbf{u}}_i^\top\hat{\mathbf{u}}_j)
       =T_c-\|\mathbf{s}_c\|_2^2+V_c.
  \label{eq:v4-category-sum}
\end{equation}
Retaining $V_c$ makes the expression exact even when numerical normalization
does not return exactly unit norm. For the background, define
\begin{equation}
  A_c^{\mathrm{out}}=\sum_{j\in B_c}\bar w_j^f,\qquad
  \mathbf{s}_c^{\mathrm{out}}=\sum_{j\in B_c}\bar w_j^f\hat{\mathbf{u}}_j.
\end{equation}
Then $U_c=A_c A_c^{\mathrm{out}}$ and
\begin{equation}
  \sum_{i\in I_c}\sum_{j\in B_c}
      w_{ic}\bar w_j^f(1-\hat{\mathbf{u}}_i^\top\hat{\mathbf{u}}_j)
       =U_c-\mathbf{s}_c^\top \mathbf{s}_c^{\mathrm{out}}.
  \label{eq:v4-background-sum}
\end{equation}
Background sums are obtained from the total family-annotated sum after
subtracting the category members' contributions using their \emph{family}
maximum confidence. This is important when a member's confidence for category
$c$ differs from its largest confidence in the family.

The implementation uses indexed reductions to evaluate these identities.
If $M_f$ is the number of annotated memberships in a family and $C_f$ its
number of categories, their additional time is proportional to
$(N_X+M_f+C_f)d_{\mathrm{bio}}$, with $d_{\mathrm{bio}}=512$; it does not require a
dense $N_X\times N_X$ biological distance matrix. This saving does not remove
the separate $N_R\times N_E$ matrix used by the full-graph association loss.

\section{Why the two association losses differ}
\label{app:v4-loss-behavior}
For a single phase-1 positive, write
$D_{ap}=\exp(s_{ap}^{(1)})+\sum_{n\in N_B(a)}\exp(s_{an}^{(1)})$. Its gradients are
\begin{align}
  \frac{\partial\ell_{\mathrm{dec}}(a,p)}{\partial s_{ap}^{(1)}}
      &=\frac{\exp(s_{ap}^{(1)})}{D_{ap}}-1,\\
  \frac{\partial\ell_{\mathrm{dec}}(a,p)}{\partial s_{an}^{(1)}}
      &=\frac{\exp(s_{an}^{(1)})}{D_{ap}},\qquad n\in N_B(a).
\end{align}
Other known positives are absent from this positive's denominator and have
zero derivative in this individual term. They receive their own positive
terms. Averaging these terms within an anchor ensures that adding more
recorded positives does not simply multiply its loss weight.

Phase 2 instead fits a categorical distribution with target probability
$1/d_a$ for each of an anchor's $d_a$ positives. Its per-logit derivative is
the predicted probability minus the target probability. Its minimum
cross-entropy for unconstrained logits approaches $\log d_a$, the entropy
of the uniform-positive target, rather than zero when $d_a>1$.
Consequently, numerical loss values from the two phases are not directly
comparable. Both losses use unobserved associations as contrastive
alternatives; neither resolves all missing-positive annotation.

\section{Ablation annotation provenance and benchmark-specific inputs}
\label{app:v4-annotations}
The biological ablation builds categories on the training endpoints of each
target. These annotations are not required to train the proposed method.
For EnzymeMap, EC assignments and atom mapping are read from the native
EnzymeMap cache only at indices recorded in the official training
associations. Mechanism labels are derived from the native mapped reactants
and products. The native quality field must be at least 0.5, and the retained
confidence is capped at one. Cofactor presence is recognized from the
training reaction participants and assigned confidence 0.4.

For ReactZyme, native EC annotations are restricted to training protein
sequences. Training reaction EC annotations use reaction-only Rhea matching
and exact canonical-chemistry lookup. Coarse bond-change categories use
the cached mapping of each training reaction, retain mapping confidence at
least 0.5, and multiply this confidence by the reaction-matching confidence.
Cofactor categories use the permitted training participant sets. These
reaction-only annotation resources provide information beyond the raw
released participant-set input; they are auxiliary training supervision,
not an inference-time reconstruction of missing reaction direction.

For every training edge, the reaction's observed categories are added to its
associated protein, retaining the maximum confidence when a membership is
supported more than once. No held-out reaction--enzyme association is used
for this propagation. Unknown or absent categories are omitted. In the
current EnzymeMap annotation vocabulary, the observed coarse mechanism groups
are phosphate transfer, redox/carbonyl interconversion, stereochemical
rearrangement, and sulfur/thiol chemistry. These broad labels should not be
described as a detailed model of catalytic mechanisms.

The target training graph and dictionary sizes are given in
Table~\ref{tab:v4-targets}. Duplicate association rows are removed when
building the compact phase-2 graph. Training endpoints can remain eligible
candidates under a benchmark's candidate-pool definition; the method does
not silently remove them.

\begin{table}[ht]
\centering
\small
\begin{tabular}{lrrr}
\toprule
Target & Training rows & Unique phase-2 edges & Reaction anchors\\
\midrule
ReactZyme Reaction-Sim & 147,393 & 147,393 & 6,977\\
ReactZyme Enzyme-Sim & 152,748 & 152,748 & 7,386\\
ReactZyme Time & 149,554 & 149,554 & 7,281\\
EnzymeMap & 34,427 & 34,180 & 12,603\\
\bottomrule
\end{tabular}
\caption{Target-specific training inputs for the implemented recipe. Each row
has independently fitted retrieval weights, phase-2 heads, and a dictionary.}
\label{tab:v4-targets}
\end{table}

\section{Optimization and inference procedure}
\label{app:v4-recipe}
\begin{longtable}{p{0.50\textwidth}p{0.42\textwidth}}
\toprule
Setting & Value\\
\midrule
\endhead
Phase-1 seed; association batch size & 42; 512\\
Retained phase-1 snapshot & End of epoch 10\\
Phase-1 optimizer & AdamW, $\mathrm{lr}=10^{-4}$, weight decay $10^{-2}$\\
Phase-1 temperature & Fixed $\tau_1=0.2$ (inverse 5)\\
Phase-1 directional weights & $1/2$ for each direction\\
Unknown-negative weight & 1\\
Entropy and diversity coefficients & 0.01 and 0.01\\
Protein attention uniform mixture & $\eta=0.05$\\
Gate uniform mixture & $\zeta=0.05$\\
Phase-2 optimizer & AdamW, $\mathrm{lr}=10^{-4}$, weight decay $10^{-3}$\\
Phase-2 updates; seed & 100 full-graph updates; 42\\
Phase-2 temperature; identity coefficient & $\tau_2=0.2$; 10\\
Phase-2 residual training coefficient & $\kappa=0.2$\\
Phase-2 gradient clipping & Global gradient norm 1\\
Added biological-label supervision & None in either training stage\\
Inference protein fusion multiplier & $\nu=3$\\
Inference residual multiplier & $c=0.5$\\
Dictionary score coefficient & $\alpha=0.4$\\
Protein dictionary neighbors & 32\\
Dictionary kernel temperatures & $\tau_E=\tau_R=0.03$\\
Reaction dictionary truncation & None; all training reaction anchors\\
\bottomrule
\end{longtable}

\paragraph{Training and inference procedure.}
\begin{enumerate}[leftmargin=*,itemsep=5pt]
  \item Construct the target training association graph and the frozen feature
  caches. Initialize the trainable retrieval components for that target and
  load the frozen pretrained feature resources and SLEEC scorer.
  \item Train the phase-1 encoders with Eq.~\eqref{eq:v4-phase1-total}. For
  each batch, deduplicate endpoint IDs and recover all known in-batch training
  positives. Retain the epoch-10 state.
  \item Freeze F3. Export native training embeddings using its learned scale
  $s$, without the factor-three inference adjustment. Build the dictionary
  centers, training protein bank, training reaction bank, and association map.
  \item Initialize the residual heads at identity. For each of 100 updates,
  encode every training endpoint, evaluate the full bidirectional
  cross-entropy and identity penalty in Eq.~\eqref{eq:v4-phase2-total}.
  Clip gradients and update only the residual heads.
  \item At inference, obtain the protein's global and fused branches and use
  $\nu s$ in Eq.~\eqref{eq:v4-inference-fusion}. Apply the trained heads with
  residual coefficient $c\kappa=0.1$.
  \item Independently compute the enzyme and reaction dictionary coordinates
  from frozen raw features and the stored training dictionary. Concatenate
  the weighted blocks, or retain equivalent dense and sparse blocks.
  \item Rank candidates by Eq.~\eqref{eq:v4-final-score} in the requested
  direction. No candidate labels, candidate-pool normalization, or biological
  annotation files are needed for encoding new endpoints.
\end{enumerate}

\paragraph{Checkpoint representation of the fusion adjustment.}
The EnzymeMap deployment checkpoint already stores the factor-three change
to the bounded fusion scalar, so its runtime fusion multiplier is one.
The ReactZyme deployment checkpoints store the native scalar and apply three
at runtime. These represent the same recipe. Applying three again to the
calibrated EnzymeMap checkpoint would implement a different model. Phase-2
training used the native feature exports for every target.

\paragraph{Numerical evaluation.}
Dictionary feature norms, centers, and nearest-neighbor similarities use
FP64 reductions where implemented before storing FP32 features. Neighbor
selection resolves cutoff ties by the lower training-anchor index.
Equivalent mathematical score decompositions can nevertheless have different
floating-point rounding when executed in different orders. The frozen
evaluation routines and candidate ordering should therefore be retained,
particularly for nearly tied ReactZyme reaction candidates. The dictionary
is part of the deployed method, not a post-hoc normalization fitted to the
current screening library.

\section{Scope of the biological-supervision ablation}
\label{app:v4-scope}
The proposed method omits biological-label supervision. The \vfourBio\
ablation adds it in the second stage; its gradients update only the residual
heads over the complete training graph,
while the residue queries and reaction encoder retain their first-stage
parameters. Category eligibility and background coverage are therefore
defined over training endpoints rather than individual minibatches.

Matched removal and shuffled-label controls show mixed, mostly small
benchmark changes. The extra unique Case1 catalyst crosses from rank 26
to 25, with unchanged recovery in the original-entry panel. These results
do not establish a consistent generalization benefit, motivating omission
of the added loss from the proposed method. Equal family coefficients do
not imply equal gradient influence or equal predictive benefit.
The 512 neural coordinates have no fixed
biochemical names; only the dictionary coordinates have an explicit
training-reaction interpretation.
